% SPDX-FileCopyrightText: Copyright (c) 2016-2026 Objectionary.com
% SPDX-License-Identifier: MIT

\section{Related Work}\label{sec:related}

In this section we analyze and categorize prior art related to our work.
Neither object-oriented formalism nor pure object-oriented languages
  are new research topics.
However, we identified certain gaps in existing studies
  that make us believe that our work has novelty.

Attempts were made to formalize OOP and introduce an object calculus
  similar to the lambda calculus~\citep{barendregt2012}
  used in functional programming.
For example, \citet{abadi1995imperative} suggested an imperative
  calculus of objects, which was extended by~\citet{bono1998imperative}
  to support classes, by~\citet{gordon1998concurrent} to support concurrency
  and synchronization, and by~\citet{jeffrey1999distributed}
  to support distributed programming.
These calculi, together with their functional counterpart,
  were consolidated into a comprehensive
  theory of objects~\citep{abadi1996theory}.

A parallel line of work embeds objects directly into the
  \(\lambda\)-calculus rather than defining a standalone calculus.
\citet{mitchell1993lambda} introduced a lambda calculus of objects
  with method specialization, extended by~\citet{di1998lambda}
  with self-inflicted extension.
These systems show that object features can be encoded on a functional core,
  which qualifies the recurring claim that OOP lacks
  a \(\lambda\)-calculus of its own; \phic{} takes the opposite route,
  treating the object rather than the function as its single primitive,
  so that dispatch and decoration are native constructs rather than encodings.

Earlier, \citet{honda1991object} combined OOP and \(\pi\)-calculus in order
  to introduce object calculus for asynchronous communication,
  which was further referenced by~\citet{jones1993pi} in their work
  on object-based design notation.

A few attempts were made to reduce existing OOP languages
  and formalize what is left.
Featherweight Java is the most notable example proposed
  by~\citet{igarashi2001featherweight}, which omits almost all features
  of the full language (including interfaces and even assignment)
  to obtain a small calculus.
Later it was extended by~\citet{jagannathan2005transactional}
  to support nested and multi-threaded transactions.
Featherweight Java is used in formal languages such as
  Obsidian~\citep{coblenz2019} and SJF~\citep{usov2020}.

The calculi above are class-based, whereas \phic{} is prototype-based:
  objects arise by copying and decorating other objects rather
  than by instantiating classes.
This tradition begins with Self~\citep{ungar1987self},
  whose delegation mechanism \citet{stein1987delegation} showed to subsume
  inheritance, and continues in calculi of delegation
  by object composition~\citep{bettinibono2008delegation}.
The meaning of inheritance in such settings has been given
  denotationally by~\citet{cook1989denotational} as a fixed point over records;
  \phic{} instead renders it operationally, through the $@$-decoration that
  its morphing operator resolves (\cref{sec:morphing}).

Another example is Larch/C++~\citep{cheon1994quick},
  which is a formal algebraic interface specification language tailored to C++.
It allows interfaces of C++ classes and functions to be documented
  in a way that is unambiguous and concise.

Several attempts to formalize OOP were made by extensions of
  the most popular formal notations and methods,
  such as Object-Z~\citep{duke1991object} and VDM++~\citep{durr1992vdm}.
In Object-Z, state and operation schemes are encapsulated into classes.
The formal model is based upon the idea
  of a class history~\citep{duke1990towards}.
However, all these OO extensions do not have comprehensive refinement
  rules that can be used to transform specifications into implemented code
  in an actual OO programming language, as was noted by~\citet{paige1999object}.

\citet{bancilhon1985calculus} suggested an object calculus
  as an extension to relational calculus.
\citet{jankowska2003anotheroop} further developed these ideas and related
  them to a Boolean algebra.
\citet{leekwakryu1996transform} developed an algorithm to transform
  an object calculus into an object algebra.

However, all these theoretical attempts to formalize OO languages were unable
  to fully describe their features,
  as was noted by~\citet{nierstrasz1991towards}: ``The development of
  concurrent object-based programming languages has suffered from
  the lack of any generally accepted formal foundations
  for defining their semantics.''
The visual and architectural notations often used to model
  object-oriented designs fare worse still:
  \citet{eden2002visual}, surveying them, observed that
  ``Not one of the notations is defined formally,
  nor provided with \nospell{denotational} semantics,
  nor founded on axiomatic semantics.''
Moreover, \citet{ciaffaglione2003reasoning,ciaffaglione2003typetheories,ciaffaglione2007theory_of_contexts}
  noted in their series of works that, at the time,
  relatively little formal work had been carried out on object-based languages.
Mechanized treatments have appeared since---among them
  a machine-checked model of a Java-like language,
  virtual machine, and compiler~\citep{klein2006jinja}---yet a compact
  object calculus whose own meta-theory is machine-checked remains uncommon.

What sets \phic{} apart from the surveyed formalisms is therefore
  not broader feature coverage, which it deliberately forgoes
  (\cref{sec:future}), but this uniformity together with
  a confluence result certified in Lean~4 (\cref{thm:normalization-confluent}).
